\chapter{Submission below Metal}\label{ch:submission-below-metal}

A Metal application asks Metal to build a pipeline and encode a command buffer; Metal's driver library
(AGXMetal) then writes command pages and calls Apple's private \code{IOGPU} framework, which enters the kernel driver. A
below-Metal process does the same work itself: it allocates GPU memory and shared memory through the AGX driver's user
client, writes the code, the launch state and the command pages, builds the submission record, and calls the same
\code{IOGPU} submission function (\cref{tab:apple-layers}); \cref{fig:submit-sequence} is that sequence. The resident runtime of \cref{ch:below-metal-native}
assembles these steps into a model execution.

\begin{xltabular}{\linewidth}{@{}>{\hsize=0.59\hsize}X>{\hsize=1.26\hsize}X>{\hsize=1.15\hsize}X@{}}
\caption{What the below-Metal process writes and what remains Apple's, by layer.}\label{tab:apple-layers}\\
\toprule
Layer & Below Metal & Still Apple's \\
\midrule
\endfirsthead
\tablecontinued{3}
\toprule
Layer & Below Metal & Still Apple's \\
\midrule
\endhead
program & every instruction byte, from the project's compiler & encoding rules recovered with Apple's decoder \\*
resources and launch & allocations, command pages, launch words, bindings, the Submit record,
written by the project's host code & the formats themselves, recovered by measurement \\*
submission & the call into \code{IOGPUCommandQueueSubmitCommandBuffers} & \code{IOGPU.framework}, IOKit, the AGX kernel driver \\*
execution &  & the GPU firmware and the hardware \\*
\bottomrule
\end{xltabular}

\leadhead{Notation.}

\begin{itemize}
\item \emph{User client}: the IOKit connection to the AGX kernel driver. \emph{Selectors}: its numbered entry points.
\item \emph{Measured graph}: a complete set of allocations, pages and launch words whose layout was captured from Metal
  once and then rebuilt from fields. Within one, the words a launch needs are known by role, except a few known only by
  value (the shader-state bits and the kernel-page word at offset 0x3c0), which are copied as captured. The graphs in
  use are listed in \cref{tab:measured-graphs}.
\item \emph{Allocation n}: the n-th GPU memory allocation in a measured graph.
\item \emph{Command pages}: the two 16 KiB shared-memory regions of a submission, the kernel page and the segment page.
\item \emph{Launch state}: the per-launch words in measured allocations that select a program, its grid and its
  buffers.
\item \emph{Submit record}: the record handed to \code{IOGPUCommandQueueSubmitCommandBuffers}.
\end{itemize}

Each claim in this chapter is tagged measured (executed below Metal with a checked output and a control), observed
(traced in a live process with Frida, or read from Apple's binary without executing a GPU command) or static (resolved
by reading the installed \code{IOGPU} code with \code{tools/g17submitinspect.c}). \code{docs/g17-first-dispatch-trials.md}
is cited as \emph{trials}, with a section title.

% The below-Metal submission sequence: what the runner writes, in order, and the
% driver entry points each step reaches. Red: the project's native runtime;
% gray: Apple's.
\begin{figure}[htbp]
\centering
\begin{tikzpicture}[node distance=3.5mm and 7mm,
    stage/.style={nat, text width=33mm, align=left, minimum height=8.5mm},
    detail/.style={tag, text width=86mm, align=left, anchor=west, font=\scriptsize}]

  \node[stage] (dev) {\textbf{1} Device setup};
  \node[detail, right=of dev] (devd)
    {AGX user client opens with no entitlement; no \code{MTLDevice} and no AGXMetal in the
     process. Queue creation calls selectors 7, 16 and 28.};

  \node[stage, below=of dev] (alloc) {\textbf{2} Allocations};
  \node[detail, right=of alloc] (allocd)
    {Selector 9 returns a buffer mapped at the same CPU and GPU address. Selector 14 gives the
     two 16\,KiB shared regions: the kernel page and the segment page.};

  \node[stage, below=of alloc] (code) {\textbf{3} Code placement};
  \node[detail, right=of code] (coded)
    {Program bytes upload into allocation 0 at 0x10000000000, from offset 0x6c0, 64-byte
     aligned; the allocation is immutable after upload.};

  \node[stage, below=of code] (launch) {\textbf{4} Launch state};
  \node[detail, right=of launch] (launchd)
    {Program counter written in pieces into allocation 23 (0x40, 0x46, 0x48), shader state at
     0x42 separately; grid, threadgroup size and buffer pointers into allocations 22, 25
     and 28.};

  \node[stage, below=of launch] (submit) {\textbf{5} Submit};
  \node[detail, right=of submit] (submitd)
    {A 48-byte record naming both page IDs and two fresh callback blocks; consecutive trace IDs;
     the fire value 0x800000f0 at offset 0x24 of shared region 0. One record goes through
     \code{IOConnectTrap4}, several through \code{IOConnectCallMethod} selector 0x1d.};

  \node[stage, below=of submit] (done) {\textbf{6} Completion};
  \node[detail, right=of done] (doned)
    {A 40-byte notification is dequeued and the referenced block invoked: one scheduled, one
     completed, in that order.};

  \foreach \a/\b in {dev/alloc, alloc/code, code/launch, launch/submit, submit/done}
    \draw[flow] (\a) -- (\b);

  \begin{scope}[on background layer]
    \node[draw=native, dashed, rounded corners=3pt, fill=native!4, inner sep=2.5mm,
          fit=(dev) (done) (devd) (submitd)] (lane) {};
  \end{scope}

  \node[band, minimum width=132mm, anchor=north] at ([yshift=-8mm]lane.south)
    (iogpu) {Apple private \code{IOGPU} framework and IOKit};
  \node[band, below=3mm of iogpu, minimum width=132mm] (drv) {Apple AGX kernel driver};
  \node[band, below=3mm of drv, minimum width=132mm] (fw) {Apple GPU firmware};

  \draw[->, line width=1.4pt, color=native] (done.south) -- (done.south |- lane.south)
    -- (done.south |- iogpu.north);
  \draw[flow] (iogpu) -- (drv);
  \draw[flow] (drv) -- (fw);

\end{tikzpicture}
\caption{The below-Metal submission sequence. Every step inside the dashed box is written by the project's
runner; the bands below it stay Apple's.}
\label{fig:submit-sequence}
\end{figure}


\section{Device setup}\label{sec:opening-device-contexts}

The process never creates an \code{MTLDevice} or a Metal command buffer, and it checks that AGXMetal is not loaded
(MM~\mm{25.210.9}).

\begin{itemize}
\item The AGX user client opens with no entitlement. A process with no Metal driver loaded creates the device and a
  context (measured, \code{tools/g17dispatch.c}, README rows 1–2).
\item Creating a command queue calls selectors 7, 16 and 28, in \code{IOGPUCommandQueueCreate} (observed in a Metal process that
  created only a device and a queue); MM~\mm{25.143} had named selector 7 the submit and fence. Selector 28 registers the
  queue's completion notification queue (static). Calling the group a second time by hand ends in a selector-28 error
  (trials, "Submission transport and record, resolved statically" and the preflight that follows it).
\item The sampler-source query (selector 261) returns a privilege error; nothing below depends on it (MM~\mm{25.143}).
\end{itemize}

\section{Allocations}\label{sec:memory-allocations-shared}

GPU allocations come from selector 9. Each request names an allocation class (a request \emph{type}), a size and an
address, and the driver returns a buffer mapped at the same CPU and GPU virtual address (measured, MM~\mm{0.15}, README
row 2). The address window depends on the request type. The exact type value 0x58000000 places an allocation in the
window at 0x6f\_00000000, while other tested values do not reach the 0x1ea\_ window that Metal's command buffer names;
whether that window is the shader window is not shown (measured, MM~\mm{25.143.O}). Allocations of 32 KiB, 8 MiB, 96 MiB
and 1 GiB were created and CPU-mapped with sentinels, but with no Submit their GPU visibility was not established, and
the runtime reuses measured graphs instead of a size formula (MM~\mm{25.210.3}).

Shared memory comes from selector 14 (\code{IOGPUDeviceCreateDeviceShmem}), which MM~\mm{25.143} had read as a telemetry
ring. Its 16-byte output is \code{\{cpu\_address: u64, size: u32, shmem\_id: u32\}}. A submission uses two 16 KiB
regions, the kernel page and the segment page, with IDs such as 2 and 1 (static and observed; trials, "Submission
transport and record").

\section{Code placement}\label{sec:code-placement-program}

Every program is uploaded into allocation 0, at GPU address 0x10000000000, and the allocation is treated as immutable
after upload. Programs start at offset 0x6c0 and are 64-byte aligned (measured, MM~\mm{25.210.6}). A launch's program
counter is allocation 0's base plus the program's offset, written in pieces into allocation 23: the low 16 bits OR 7 at
offset 0x40, bits 16–31 at offset 0x46 and bits 32–47 at offset 0x48 (MM~\mm{25.210.6}). The program counter's low tag bits are known
only for the one entry class these programs use. The shader-state word at offset 0x42 of
allocation 23 is written separately, because a single 32-bit program-counter write once corrupted it (measured,
MM~\mm{25.210.6}).

\section{Launch state}\label{sec:launch-state-command}

The runner writes the program counter, the shader-state and scratch words, the grid, the threadgroup size, up to three buffer pointers and two
scaled address coordinates into allocations 22, 23, 25 and 28, and keeps three kernel-page words as captured
(\cref{tab:launch-words}).

\begin{xltabular}{\linewidth}{@{}>{\hsize=1.14\hsize}X>{\hsize=0.67\hsize}X>{\hsize=1.38\hsize}X>{\hsize=0.81\hsize}X@{}}
\caption{Launch words, all written by the project's runner (MM~\mm{25.210.6} unless noted).}\label{tab:launch-words}\\
\toprule
Where & Width & Meaning & Status \\
\midrule
\endfirsthead
\tablecontinued{4}
\toprule
Where & Width & Meaning & Status \\
\midrule
\endhead
allocation 23, offsets 0x40, 0x46, 0x48 & 16 bits each & program counter: low half OR 7, bits 16–31, bits 32–47 & measured \\
allocation 23, offset 0x42 & 16 bits & shader state & measured value; bits unnamed \\
allocation 23, offset 0x38 & two 32-bit words & scratch/config word, then zero & measured value; bits unnamed \\
allocation 22 offset 0xa8; allocation 25 offset 0x10 & three 32-bit words each & grid x, y, z, written twice & measured \\
allocation 25, offset 0x1c & 32 bits & threads per threadgroup (32 in the measured classes) & measured \\
allocation 28, offset 0x1ba0 & four 64-bit pointers & program buffer pointers; at most three are admitted & measured \\
allocation 23, offset 0x06 & 16 bits & scaled metadata address coordinate: \code{(VA28 >> 13) \& 0xffff} & measured (see below) \\
allocation 25, offset 0x09 & 16 bits & scaled address coordinate: \code{(VA23 >> 14) \& 0xffff} & measured equality \\
kernel page, offset 0x3c0 & word & needed for the tensor contribution: clearing it removes the MMA result
while Submit and callbacks succeed & measured; meaning unknown \\
kernel page, offsets 0x3ec and 0x3f8 & words & kept as measured & unknown \\
\bottomrule
\end{xltabular}

Scalar and packing programs use the state word 0x0e40 with the scratch word 0x0c000007. Tensor
projections and the cooperative LayerNorm use 0x0e58 with 0x0c00100f (MM~\mm{25.210.6}).

The halfwords at allocation 23 offset 0x06 and allocation 25 offset 0x09 fit the shifted addresses above exactly at
six graph sizes from 17 to 34 MiB and in the 4 MiB graph (trials, the address-coordinate table), although
MM~\mm{25.210.6} lists them only as "graph-identity correlates; no formula". For allocation 23's halfword the evidence is
causal in one graph: changing only that halfword between the high and low coordinate decided pass or fail. No
measurement shows which firmware step reads it or whether every allocation with that encoding has the same role.

In the 1 GiB class the complete 32 KiB shader record is copied into allocation 17 at offset 0x10000 and
selected by allocation 25's halfword (0x001a) at offset 0x09; the originally wrapped selector value (0x0048)
page-faulted and must not be reused. The 48 KiB binding metadata is copied from allocation 28 into allocation 17 and
selected by allocation 23's halfword (0x002c) at offset 0x06. The executor rewrites the selected low copy of the shader
record for each stage; writing the original high record after the copy does not change the program that runs
(measured, MM~\mm{25.210.3}).

\section{The Submit call}\label{sec:submit-record-call}

The Submit record's base layout is 48 bytes; Metal's observed records are 64 bytes, one per command buffer
(\cref{tab:submit-record}). The record layout and the call paths in this section were resolved statically and observed
in Metal processes (trials, "Submission transport and record").

\begin{xltabular}{\linewidth}{@{}lX@{}}
\caption{Fields of the Submit record by byte offset.}\label{tab:submit-record}\\
\toprule
Offset & Field \\
\midrule
\endfirsthead
\tablecontinued{2}
\toprule
Offset & Field \\
\midrule
\endhead
0x00, u32 & shared-memory ID of the kernel page \\*
0x04, u32 & shared-memory ID of the segment page \\*
0x08, u32 & optional sideband shared-memory ID \\*
0x10, u64 & scheduled callback block \\*
0x18, u64 & completed callback block \\*
0x20, u32 & optional debug shared-memory ID \\*
0x28, u64 & declared by the framework's metadata; meaning not
established \\*
0x30 on & possible driver extension; kept uninterpreted \\*
\bottomrule
\end{xltabular}

A Submit also needs consecutive trace IDs, the fire value and fresh callback blocks (measured, MM~\mm{25.210.6}):

\begin{itemize}
\item \code{IOGPUDeviceGetNextGlobalTraceID} must return consecutive IDs for the trace-ID words in the command pages.
\item The coordination word at offset 0x24 of shared region 0 must hold the fire value 0x800000f0. The staged value 0xf0 was
  rejected.
\item Exactly one scheduled and one completed callback are required, each a fresh \code{Block\_copy}. Reusing blocks crashed the
  CPU notification dispatcher (a host-side crash).
\end{itemize}

A one-record and a multi-record Submit reach the kernel through different calls (static):

\begin{itemize}
\item A one-record Submit calls \code{IOConnectTrap4}: a kernel trap with trap index 0, the queue ID, the record size, the record
  pointer and an output pointer. The Metal reference runs observed exactly this, with count 1 and stride 64.
\item A multi-record Submit calls \code{IOConnectCallMethod} with selector 0x1d, passing the contiguous record array.
\end{itemize}

MM~\mm{25.143} and README row 6a, which describe a Submit with no per-dispatch kernel call, name a candidate store at
Submit+0x11c. That store is an output-argument write (\cref{app:a}).

In a Metal process, a control that never set the ready bit still executed (MM~\mm{25.143}). That observation does not show what
consumes the fire value, or whether writing 0x800000f0 before Submit, as the runner does, is what starts the GPU.

\section{Completion}\label{sec:completion-callbacks-their}

Completion arrives as a notification. \code{IOGPUNotificationQueueDispatchAvailableCompletionNotifications} dequeues
40-byte notifications (\code{IODataQueueDequeue}) and invokes the referenced callback block. In the Metal reference runs,
two notifications per submission carried, as their first 64-bit values, exactly the record's scheduled and completed
callback pointers, in that order (observed; trials, "Submission transport and record"). In the MiniLM FFN block, a warm
request spends 2.385~ms from Submit to the completed-callback entry and 0.0495~ms from callback entry to the waiter's
return (measured, MM~\mm{25.210.5}). The first interval includes driver work and completion delivery, so it overstates
GPU time.

Callback ownership was observed on macOS 26.6.2 (MM~\mm{25.210.9a}). On the successful count-one path, each of six
observed callback blocks was invoked, returned, released by the notification dispatcher, and freed through
\code{\_Block\_release}. No caller release occurred, and no block accumulated. Ownership on failure, timeout or
queue teardown with pending callbacks, and on other OS versions, is not known.

\section{Multi-record Submits and dependent dispatches}\label{sec:several-records-per}

Several records in one Submit execute. Count-two and count-three Submits ran distinct programs and bindings into
separate outputs, and later trials reached seven separately bound records, the seventh predicted without a Metal
capture (measured; trials, "Two command records in one pure Submit" through "Seventh record predicted").

Dependent stages, as the resident runtime runs them, are consecutive count-one Submits: 7 for the MiniLM FFN, 916 for
one Qwen decode token (measured, MM~\mm{25.210.3}, MM~\mm{25.210.5}). A two-stage dependent pipeline, with and without
a program switch, passed in the trials (measured; trials, "Dependent two-stage pure GPU pipeline"). The runtime does
not issue many dependent dispatches in one Submit with explicit barriers between them, although Metal's
serial-barrier word and fence semantics are decoded (README row 10).

\section{Faults}\label{sec:faults-hangs-recovery}

Malformed kernels have hung the GPU on this machine. Killing the host process does not cancel a submitted command
buffer, and recovery from that hang was a reboot (rule 35 of MM~\mm{17}). The GPU's recovery counter and event reports
are observable before and after a run, and the resident runtime refuses a run if either changes (MM~\mm{25.210.9}). It
bounds the wait for each Submit's callbacks to five seconds (MM~\mm{25.210.6}), and it cannot detect or recover from a
hang beyond that bound or reset the GPU (README row 11).

\section{Scope of the measured graphs}\label{sec:measured-graphs-generalizes}

The runtime uses four measured graphs (\cref{tab:measured-graphs}) and refuses every other size.

\begin{xltabular}{\linewidth}{@{}l>{\hsize=1.27\hsize}X>{\hsize=0.73\hsize}X@{}}
\caption{Measured resource graphs.}\label{tab:measured-graphs}\\
\toprule
Graph & Used for & Source \\
\midrule
\endfirsthead
\tablecontinued{3}
\toprule
Graph & Used for & Source \\
\midrule
\endhead
3 MiB & the MiniLM FFN and attention blocks & MM~\mm{25.210.4}, MM~\mm{25.210.6} \\*
4 MiB & the connected layer (allocation 20 enlarged to 4 MiB, allocation 21
moved) & MM~\mm{25.210.6} \\*
128 MiB & the MiniLM encoder & MM~\mm{25.210.3} \\*
1 GiB & the Qwen2.5-0.5B decoder & MM~\mm{25.210.3} \\*
\bottomrule
\end{xltabular}

Outside these graphs the program-counter placement, the grid and threadgroup words, up to three buffer pointers, the
two state classes, the record layout, and the address-coordinate halfwords within the tested range still apply. These
have not been shown to work:

\begin{itemize}
\item arbitrary graph allocation and publication (the runtime recovers only the request and page sets of the measured
  graphs);
\item a fourth buffer pointer (earlier controls wrote 0 of 32 outputs);
\item new shapes, data types, simdgroup counts, OS versions or GPU generations;
\item GPU-only timestamps.
\end{itemize}

\textbf{Evidence.} MM~\mm{0.15}, MM~\mm{25.143} (with 25.143.N and 25.143.O), MM~\mm{25.147}, MM~\mm{25.150}, MM~\mm{25.210} (25.210.3, 25.210.5,
25.210.6, 25.210.9, 25.210.9a); \code{docs/g17-first-dispatch-trials.md}; README rows 1–16.
